\documentclass[reqno]{amsart}
\usepackage[T1]{fontenc}
\usepackage{amsmath,amssymb,amsthm}
\usepackage{xr-hyper}
\usepackage[hidelinks]{hyperref}
\externaldocument[sd-]{20261125T000000--slope-method__workbook}
\externaldocument[at-]{20261126T030100--a-twist-for-each-prime__workbook}
\externaldocument[da-]{20261126T030300--the-data__workbook}
\externaldocument[ta-]{20261126T030600--theorem-a-short-forms__workbook}
\externaldocument[dch-]{20261126T030700--the-determinant-of-a-chain__workbook}
\externaldocument[col-]{20261126T030800--the-cost-of-one-lift__workbook}
\externaldocument[tw-]{20261126T031000--twists-with-few-congruences__workbook}
\externaldocument[mi-]{20261126T031100--theorem-c-and-the-main-inequality__workbook}
\externaldocument[ob-]{20261126T040100--spaces-forms-and-twists__workbook}

% Chapter number; \ref{chap} gives it.
\newcounter{chap}
% Chapter numbers, generated from 20261125T012000--grh-workbook__guide.org
% by scripts/make-chapters.py.  Do not edit: change the order of the
% index and rerun the script (./preview-site.sh and CI do this).
%
% Usage in a sheet, after \newcounter{chap}:
%   % Chapter numbers, generated from 20261125T012000--grh-workbook__guide.org
% by scripts/make-chapters.py.  Do not edit: change the order of the
% index and rerun the script (./preview-site.sh and CI do this).
%
% Usage in a sheet, after \newcounter{chap}:
%   % Chapter numbers, generated from 20261125T012000--grh-workbook__guide.org
% by scripts/make-chapters.py.  Do not edit: change the order of the
% index and rerun the script (./preview-site.sh and CI do this).
%
% Usage in a sheet, after \newcounter{chap}:
%   \input{chapters}\chapnumber{<sheet ID>}
\makeatletter
\providecommand\chapters@set[2]{\@namedef{chapters@#1}{#2}}
\chapters@set{20261125T001000}{1}% zeros-to-ramp-sums
\chapters@set{20261125T001100}{2}% inert-primes-and-estimate-b
\chapters@set{20261125T005000}{3}% galois-fourier-and-averaging
\chapters@set{20261125T005500}{4}% columns-from-a-sphere
\chapters@set{20261125T000000}{5}% slope-method
\chapters@set{20261126T030100}{6}% a-twist-for-each-prime
\chapters@set{20261125T003000}{7}% the-chain-and-the-square-root
\chapters@set{20261125T010000}{8}% rank-lemma-and-zero-estimate
\chapters@set{20261125T010500}{9}% poincare-series-and-near-orthogonality
\chapters@set{20261124T223000}{10}% elliptic-curves-for-the-proof
\chapters@set{20261125T009500}{11}% fock-model-and-theta-intertwiners
\chapters@set{20261124T223100}{12}% complex-multiplication-and-inert-reduction
\chapters@set{20261125T008000}{13}% one-cm-determinant
\chapters@set{20261124T214500}{14}% supersingular-theta-divisibility
\chapters@set{20261124T214600}{15}% the-isotropic-plane
\chapters@set{20261125T013700}{16}% pairings-on-a-formal-group-of-height-two
\chapters@set{20261125T013500}{17}% hecke-nondegeneracy-on-the-tree
\chapters@set{20261126T030200}{18}% the-bound-is-attained
\chapters@set{20261125T009000}{19}% kronecker-limit-and-chowla-selberg
\chapters@set{20261125T002000}{20}% definite-theta-lifts
\chapters@set{20261125T015500}{21}% siegel-weil-in-the-smallest-case
\chapters@set{20261125T002100}{22}% definite-theta-lifts-and-norms
\chapters@set{20261125T016500}{23}% hecke-operators-and-theta-lifts
\chapters@set{20261125T016700}{24}% divisibility-along-the-chain
\chapters@set{20261125T004000}{25}% denominators-and-sha
\chapters@set{20261125T004100}{26}% tate-twists-and-sha
\chapters@set{20261126T020430}{27}% the-minus-part-by-gauss-sums
\chapters@set{20261125T014100}{28}% congruence-modules
\chapters@set{20261125T014500}{29}% eisenstein-congruences-and-selmer-classes
\chapters@set{20261125T014700}{30}% pseudocharacters-and-the-length-bound
\chapters@set{20261125T011000}{31}% galois-representations-black-box
\chapters@set{20261126T030300}{32}% the-data
\chapters@set{20261126T020210}{33}% bessel-poincare-series-and-their-gram-operators
\chapters@set{20261126T030400}{34}% the-seesaw-identity
\chapters@set{20261126T030500}{35}% divisibility-of-the-composite-lift
\chapters@set{20261126T030600}{36}% theorem-a-short-forms
\chapters@set{20261126T030700}{37}% the-determinant-of-a-chain
\chapters@set{20261126T030800}{38}% the-cost-of-one-lift
\chapters@set{20261126T030900}{39}% denominators-of-projections-to-lifts
\chapters@set{20261126T031000}{40}% twists-with-few-congruences
\chapters@set{20261126T031100}{41}% theorem-c-and-the-main-inequality
\providecommand\chapnumber[1]{%
  \@ifundefined{chapters@#1}%
    {\PackageError{chapters}{Sheet #1 is not listed in the index}%
      {List it in 20261125T012000--grh-workbook__guide.org\MessageBreak
       and run scripts/make-chapters.py.}}%
    {\edef\thechap{\@nameuse{chapters@#1}}}}
\makeatother
\chapnumber{<sheet ID>}
\makeatletter
\providecommand\chapters@set[2]{\@namedef{chapters@#1}{#2}}
\chapters@set{20261125T001000}{1}% zeros-to-ramp-sums
\chapters@set{20261125T001100}{2}% inert-primes-and-estimate-b
\chapters@set{20261125T005000}{3}% galois-fourier-and-averaging
\chapters@set{20261125T005500}{4}% columns-from-a-sphere
\chapters@set{20261125T000000}{5}% slope-method
\chapters@set{20261126T030100}{6}% a-twist-for-each-prime
\chapters@set{20261125T003000}{7}% the-chain-and-the-square-root
\chapters@set{20261125T010000}{8}% rank-lemma-and-zero-estimate
\chapters@set{20261125T010500}{9}% poincare-series-and-near-orthogonality
\chapters@set{20261124T223000}{10}% elliptic-curves-for-the-proof
\chapters@set{20261125T009500}{11}% fock-model-and-theta-intertwiners
\chapters@set{20261124T223100}{12}% complex-multiplication-and-inert-reduction
\chapters@set{20261125T008000}{13}% one-cm-determinant
\chapters@set{20261124T214500}{14}% supersingular-theta-divisibility
\chapters@set{20261124T214600}{15}% the-isotropic-plane
\chapters@set{20261125T013700}{16}% pairings-on-a-formal-group-of-height-two
\chapters@set{20261125T013500}{17}% hecke-nondegeneracy-on-the-tree
\chapters@set{20261126T030200}{18}% the-bound-is-attained
\chapters@set{20261125T009000}{19}% kronecker-limit-and-chowla-selberg
\chapters@set{20261125T002000}{20}% definite-theta-lifts
\chapters@set{20261125T015500}{21}% siegel-weil-in-the-smallest-case
\chapters@set{20261125T002100}{22}% definite-theta-lifts-and-norms
\chapters@set{20261125T016500}{23}% hecke-operators-and-theta-lifts
\chapters@set{20261125T016700}{24}% divisibility-along-the-chain
\chapters@set{20261125T004000}{25}% denominators-and-sha
\chapters@set{20261125T004100}{26}% tate-twists-and-sha
\chapters@set{20261126T020430}{27}% the-minus-part-by-gauss-sums
\chapters@set{20261125T014100}{28}% congruence-modules
\chapters@set{20261125T014500}{29}% eisenstein-congruences-and-selmer-classes
\chapters@set{20261125T014700}{30}% pseudocharacters-and-the-length-bound
\chapters@set{20261125T011000}{31}% galois-representations-black-box
\chapters@set{20261126T030300}{32}% the-data
\chapters@set{20261126T020210}{33}% bessel-poincare-series-and-their-gram-operators
\chapters@set{20261126T030400}{34}% the-seesaw-identity
\chapters@set{20261126T030500}{35}% divisibility-of-the-composite-lift
\chapters@set{20261126T030600}{36}% theorem-a-short-forms
\chapters@set{20261126T030700}{37}% the-determinant-of-a-chain
\chapters@set{20261126T030800}{38}% the-cost-of-one-lift
\chapters@set{20261126T030900}{39}% denominators-of-projections-to-lifts
\chapters@set{20261126T031000}{40}% twists-with-few-congruences
\chapters@set{20261126T031100}{41}% theorem-c-and-the-main-inequality
\providecommand\chapnumber[1]{%
  \@ifundefined{chapters@#1}%
    {\PackageError{chapters}{Sheet #1 is not listed in the index}%
      {List it in 20261125T012000--grh-workbook__guide.org\MessageBreak
       and run scripts/make-chapters.py.}}%
    {\edef\thechap{\@nameuse{chapters@#1}}}}
\makeatother
\chapnumber{<sheet ID>}
\makeatletter
\providecommand\chapters@set[2]{\@namedef{chapters@#1}{#2}}
\chapters@set{20261125T001000}{1}% zeros-to-ramp-sums
\chapters@set{20261125T001100}{2}% inert-primes-and-estimate-b
\chapters@set{20261125T005000}{3}% galois-fourier-and-averaging
\chapters@set{20261125T005500}{4}% columns-from-a-sphere
\chapters@set{20261125T000000}{5}% slope-method
\chapters@set{20261126T030100}{6}% a-twist-for-each-prime
\chapters@set{20261125T003000}{7}% the-chain-and-the-square-root
\chapters@set{20261125T010000}{8}% rank-lemma-and-zero-estimate
\chapters@set{20261125T010500}{9}% poincare-series-and-near-orthogonality
\chapters@set{20261124T223000}{10}% elliptic-curves-for-the-proof
\chapters@set{20261125T009500}{11}% fock-model-and-theta-intertwiners
\chapters@set{20261124T223100}{12}% complex-multiplication-and-inert-reduction
\chapters@set{20261125T008000}{13}% one-cm-determinant
\chapters@set{20261124T214500}{14}% supersingular-theta-divisibility
\chapters@set{20261124T214600}{15}% the-isotropic-plane
\chapters@set{20261125T013700}{16}% pairings-on-a-formal-group-of-height-two
\chapters@set{20261125T013500}{17}% hecke-nondegeneracy-on-the-tree
\chapters@set{20261126T030200}{18}% the-bound-is-attained
\chapters@set{20261125T009000}{19}% kronecker-limit-and-chowla-selberg
\chapters@set{20261125T002000}{20}% definite-theta-lifts
\chapters@set{20261125T015500}{21}% siegel-weil-in-the-smallest-case
\chapters@set{20261125T002100}{22}% definite-theta-lifts-and-norms
\chapters@set{20261125T016500}{23}% hecke-operators-and-theta-lifts
\chapters@set{20261125T016700}{24}% divisibility-along-the-chain
\chapters@set{20261125T004000}{25}% denominators-and-sha
\chapters@set{20261125T004100}{26}% tate-twists-and-sha
\chapters@set{20261126T020430}{27}% the-minus-part-by-gauss-sums
\chapters@set{20261125T014100}{28}% congruence-modules
\chapters@set{20261125T014500}{29}% eisenstein-congruences-and-selmer-classes
\chapters@set{20261125T014700}{30}% pseudocharacters-and-the-length-bound
\chapters@set{20261125T011000}{31}% galois-representations-black-box
\chapters@set{20261126T030300}{32}% the-data
\chapters@set{20261126T020210}{33}% bessel-poincare-series-and-their-gram-operators
\chapters@set{20261126T030400}{34}% the-seesaw-identity
\chapters@set{20261126T030500}{35}% divisibility-of-the-composite-lift
\chapters@set{20261126T030600}{36}% theorem-a-short-forms
\chapters@set{20261126T030700}{37}% the-determinant-of-a-chain
\chapters@set{20261126T030800}{38}% the-cost-of-one-lift
\chapters@set{20261126T030900}{39}% denominators-of-projections-to-lifts
\chapters@set{20261126T031000}{40}% twists-with-few-congruences
\chapters@set{20261126T031100}{41}% theorem-c-and-the-main-inequality
\providecommand\chapnumber[1]{%
  \@ifundefined{chapters@#1}%
    {\PackageError{chapters}{Sheet #1 is not listed in the index}%
      {List it in 20261125T012000--grh-workbook__guide.org\MessageBreak
       and run scripts/make-chapters.py.}}%
    {\edef\thechap{\@nameuse{chapters@#1}}}}
\makeatother

\ifcsname chapters@20261126T040200\endcsname
  \chapnumber{20261126T040200}
\else
  \renewcommand{\thechap}{?}
\fi
\renewcommand{\thesection}{\thechap.\arabic{section}}

\theoremstyle{definition}
\newtheorem{problem}{Problem}[section]
\theoremstyle{plain}
\newtheorem{theorem}{Theorem}[section]

\newcommand{\Z}{{\mathbb{Z}}}
\newcommand{\Q}{{\mathbb{Q}}}
\newcommand{\R}{\mathbb R}
\newcommand{\C}{{\mathbb{C}}}
\newcommand{\Oc}{\mathcal O}
\newcommand{\Lc}{\mathfrak L}
\newcommand{\main}{h_{\mathrm{main}}}
\newcommand{\rem}{\mathrm{rem}}
\newcommand{\hdir}{h_{\mathrm{dir}}}
\newcommand{\hgt}{\operatorname{ht}}
\newcommand{\Gram}{\operatorname{Gram}}
\newcommand{\lean}[1]{\texttt{#1}}
\setlength{\emergencystretch}{3em}

\title{Theorem C}
\author{}
\date{November 26, 2026}

\begin{document}
\maketitle
\refstepcounter{chap}\label{chap}

\paragraph{Prerequisites.}
Hadamard's and Fischer's inequalities for positive definite Hermitian matrices \cite[\S7.8]{HJ}.  From Chapter~\ref{sd-chap} heights, covolume and Problem~\ref{sd-pr:height-mul}; from Chapter~\ref{at-chap} the theorem on a twist for each prime; from Chapter~\ref{da-chap} the data and Problem~\ref{da-pr:main}; from Chapter~\ref{ob-chap} the objects and the target lattice; and the statements of the auxiliary forms, the factorization of the chain determinant and the bound for small ratios (Chapter~\ref{dch-chap}), of small ratios and the cost of one lift (Chapter~\ref{col-chap}), and of regular families (Chapter~\ref{tw-chap}).  Section~\ref{sec:toy} uses no earlier chapter.

\medskip
Theorem~C gives the lower bound used in Chapter~\ref{mi-chap}, \emph{The main inequality}, for the covolume plus twice the index of a linearly independent family of target forms spanning almost all of the target.  For each prime \(\ell\) and each field isomorphism \(\iota\colon\C\to\overline{\Q}_\ell\), the proof chooses a twist \(\theta_\iota\).  It compares the chain Gram determinant at the trivial twist with its value at \(\theta_\iota\), using a height bound for the family of ratios for each single lift.  We first compare index and Gram determinant in a lattice, then state and prove Theorem~C in the form used in the proof.  This chapter heads Part~IX; the chapters that follow it develop its inputs.

\section{The comparison in a lattice}\label{sec:toy}

Let \(N \ge 1\) be an integer, and give \(\C^d\) the Hermitian form \(\langle v, w \rangle := N\, \bar v^{\mathsf T} w\).  For \(y_1, \dots, y_n \in \Q^d\) write \(\Gram(y) := \det(\langle y_i, y_j \rangle)_{i,j \le n}\), and let \(H := \log N\).

\begin{problem}\label{pr:toy-int}
If \(y_1, \dots, y_d \in \Z^d\) are linearly independent, then \(\Gram(y) \ge N^d\), so \(\log \Gram(y) \ge dH\).  For \(N = 5\), \(y_1 = (1, 2)\), \(y_2 = (2, -1)\) one has \(\Gram(y) = 625\).
\end{problem}

\begin{problem}\label{pr:toy-index}
Let \(y_1, \dots, y_d \in \Q^d\) be linearly independent, and let \(u \in M_d(\Z)\) with \(\det u \ne 0\) be such that \(\sum_k u_{ik} y_k \in \Z^d\) for every \(i\).  If \(\log \Gram(y) \le c\) and \(\log\lvert\det u\rvert \le \delta\), then \(c + 2\delta \ge dH\).  For \(N = 5\), \(y_1 = (\tfrac12, 1)\), \(y_2 = (0, \tfrac13)\), \(u = \mathrm{diag}(2, 3)\), \(c = \log \Gram(y)\) and \(\delta = \log 6\) there is equality.
\end{problem}

Now let the family be shorter than the dimension.  Suppose it is completed to a basis by vectors of bounded length.

\begin{problem}\label{pr:toy-complete}
Let \(y_1, \dots, y_n\) and \(z_1, \dots, z_{n'}\) in \(\Z^d\) together form a basis of \(\Q^d\), with \(\langle z_i, z_i \rangle \le e^{\beta}\) for every \(i\).  Then \(\log \Gram(y) + n'\beta \ge dH\).  For \(d = 3\), \(N = 2\), \(y_1 = (1, 1, 0)\), \(y_2 = (0, 1, 1)\), \(z_1 = (1, 0, 0)\): \(\Gram(y) = 12\), \(\Gram(y, z) = 8\), and \(8 \le 12 \cdot 2\).
\end{problem}

\begin{problem}\label{pr:toy-steinitz}
Every linearly independent family \(y_1, \dots, y_n \in \Z^d\) is completed to a basis of \(\Q^d\) by \(d - n\) of the standard basis vectors.  Hence \(\log \Gram(y) \ge nH\).
\end{problem}

In Section~\ref{sec:thmC} the main amount enters through the height of the family \(\iota \mapsto Z(1)/Z(\theta_\iota)\), where the twist \(\theta_\iota\) is chosen for each pair \((\ell, \iota)\) (Problem~\ref{pr:volume}).

\section{Theorem C}\label{sec:thmC}

We use the parameters \(q_0,d\) and the data of Chapter~\ref{da-chap}, and the objects and notation of Chapter~\ref{ob-chap}: the target forms \(T\), of dimension \(d_T\); the set \(\mathrm{Tw}\) of \(m\)-fold twists \(\vartheta\), with the trivial twist \(1\) and the norm \(\mathrm{Nm}\,Q_\vartheta\) of the modulus; the lower chain \(\Theta_\vartheta\); the Gram determinants \(Z(\vartheta)\) of the chain and \(Z_j(\vartheta)\) of the \(j\)th single lift at \(\vartheta^j\); and families of twists \(\theta = (\theta_\iota)\), one twist for each \(\iota \mid \ell\).  Among the conditions on the data, the scales \(x < x_1 < b\) are integers with
\[
  (\log b)^{100} \le x , \qquad x^{10^6} \le x_1 , \qquad x_1^{10^6} \le b , \qquad (\log x_1)^6 \le \log b , \qquad \log\log b \ge 10^{12} ,
\]
the last inequality following from \(10^{12}C\le\log\log b\) and \(C\ge1\), with \(C\) as in Chapter~\ref{da-chap}, Section~\ref{da-sec:data}.  Further \(a = xb\), \(\sqrt x \le m \le 2\sqrt x\), \(x^5 \le r \le 3x^5\), \(k = r + 1 - m\), \(t = 5r\) and \(f \ge 1\).  With \(U = fa\log^3 b\) write \(F := U\log b\).  The main amount is \(\main = mk\,\hdir\), and \(\main \ge 0\) by Problem~\ref{da-pr:main}.  Heights of families \(\iota \mapsto u_\iota\) are those of Chapter~\ref{sd-chap}, Section~\ref{sd-sec:heights}; we write \(\hgt(u) \le c\).

\emph{The target lattice.}  The target lattice on \(T\), with twists \(\mathrm{Tw}\), the injective lifts \(\Theta_\vartheta\), the twists allowed at \(\ell\), the lattices \(\Lc_{\vartheta,\iota}\), the norms \(\lVert v \rVert_\sigma\), and index and covolume of a family in \(T\), is that of Definition~\ref{ob-def:target-lattice}; its volume at \(\vartheta\) is \(Z(\vartheta)\).

Theorem~\ref{at-thm:twist}, for this lattice, reads as follows.  Let \(y_1, \dots, y_n, z_1, \dots, z_{n'}\) be a basis of \(T\), where \(y\) has covolume at most \(c\) and index at most \(\delta\); let \(\theta\) be a family of twists with \(\theta_\iota\) allowed at \(\ell\), and \(\gamma \colon \mathrm{Tw} \to \C^\times\); suppose (i) \(\gamma(\vartheta)\Theta_\vartheta z_i \in \Lc_{\vartheta,\iota}\) for every \(i\), every \(\iota \mid \ell\) and every \(\vartheta\) allowed at \(\ell\); (ii) \(\lVert z_i \rVert_\sigma^2 \le e^{c_{\mathrm{short}}}\) for every \(i\) and \(\sigma\); (iii) \(\hgt(\lvert\gamma(\theta_\iota)\rvert^{-2}) \le c_{\mathrm{den}}\); (iv) \(\hgt(Z(1)/Z(\theta_\iota)) \le c_{\mathrm{vol}}\).  Then
\begin{equation}\label{eq:twist}
  0 \le c + 2\delta + n'(c_{\mathrm{short}} + c_{\mathrm{den}}) + c_{\mathrm{vol}} .
\end{equation}

Four inputs feed it.  (a) \emph{Auxiliary forms} (Theorem~\ref{dch-thm:aux}).  There are an integer \(e_\gamma\) and \(\gamma \colon \mathrm{Tw} \to \C^\times\) with \(\gamma(1) = 1\) and \(\lvert\gamma(\vartheta)\rvert = (\mathrm{Nm}\,Q_\vartheta)^{e_\gamma/2}\prod_{0 \le i, j < m} \rho'_{ij}(\vartheta)\) such that \(T\) is spanned by forms \(z\) with \(\gamma(\vartheta)\Theta_\vartheta z\) integral at \(\iota\) whenever \(\gcd(\ell, \mathrm{Nm}\,Q_\vartheta) = 1\) and with \(\lVert z \rVert_\sigma^2 \le \exp(2x_1^2b^2)\) for every \(\sigma\).  (b) \emph{Single lifts.}  There is a family \(\theta\) with small ratios (Definition~\ref{col-def:small-ratios}) such that \(\hgt(Z_j(1)/Z_j(\theta_\iota)) \le d_T(F/m - k\hdir)\) for \(1 \le j \le m\); this combines the cost of one lift (Chapter~\ref{col-chap}) with a regular family (Chapter~\ref{tw-chap}).  Each \(\theta_\iota\) of such a family is allowed at \(\ell\).  (c) \emph{Factorization} (Theorem~\ref{dch-thm:main}).  There is an integer \(e_Z\) such that for every \(\vartheta \in \mathrm{Tw}\)
\[
  \frac{Z(\vartheta)}{Z(1)} = \Bigl( (\mathrm{Nm}\,Q_\vartheta)^{e_Z/2} \prod_{1 \le i < j \le m} \rho_{ij}(\vartheta) \Bigr)^{d_T} \prod_{j=1}^m \frac{Z_j(\vartheta)}{Z_j(1)} .
\]
(d) \emph{Small ratios} (Theorem~\ref{dch-thm:small-ratios}).  If \(\theta\) has small ratios, then for every integer \(e\) both families \(\bigl((\mathrm{Nm}\,Q_{\theta_\iota})^{e/2}\prod_{0 \le i,j < m} \rho'_{ij}(\theta_\iota)\bigr)^{-1}\) and \(\bigl((\mathrm{Nm}\,Q_{\theta_\iota})^{e/2}\prod_{1 \le i < j \le m} \rho_{ij}(\theta_\iota)\bigr)^{-1}\) have height at most \(m^2tx_1\).

\begin{problem}\label{pr:scales}
The data satisfy \(2m^2tx_1 \le b\), \(4(2x_1^2 + b^{-1}) \le b\), \(b \le F\) and \(4/b^2 \le 1/2\).
\end{problem}

\begin{problem}\label{pr:completion}
If \(n \ge (1 - 4/b^2)d_T\) and \(d_T = n + n'\), then \(n'(2x_1^2b^2 + b) \le d_T b\).
\end{problem}

So the allowance for one completing auxiliary form is \(2x_1^2b^2 + b\), and the total allowance for the \(n'\) completing forms is at most \(d_T b\).

\begin{problem}\label{pr:denom}
For \(\gamma\) as in (a) and \(\theta\) as in (b), \(\hgt(\lvert\gamma(\theta_\iota)\rvert^{-2}) \le b\).
\end{problem}

\begin{problem}\label{pr:volume}
For \(\theta\) as in (b), \(\hgt(Z(1)/Z(\theta_\iota)) \le d_T(b + F - \main)\).
\end{problem}

\begin{problem}\label{pr:C}
Let \(y_1, \dots, y_n \in T\) be linearly independent, with \(n \ge (1 - 4/b^2)d_T\), covolume at most \(c\) and index at most \(\delta\).  Then \(T\) has a basis \(y_1, \dots, y_n, z_1, \dots, z_{n'}\) with each \(z_i\) as in (a), and \(c + 2\delta \ge d_T(\main - 3F)\).
\end{problem}

Problems~\ref{pr:scales}--\ref{pr:C} prove the following.

\begin{theorem}[Theorem C]\label{thm:C}
Let \(y_1, \dots, y_n\) be linearly independent target forms with \((1 - 4/b^2)d_T \le n\), of index at most \(\delta\) and covolume at most \(c\).  Then
\[
  d_T(\main - 3fa\log^4 b) \le c + 2\delta .
\]
\end{theorem}

\paragraph{In the Lean formalization.}  In the Lean proof~\cite{GL}, Theorem~\ref{thm:C} is \lean{New.Data.integral\_forms\_long}, in \lean{grh/src/New/Determinant/Twist/LongForms.lean}, and the input~(b) is \lean{New.Data.exists\_hasSmallRatios\_heightLe\_liftDeterminant\_div} in the same file; \eqref{eq:twist} is the case of the target lattice of \lean{New.TwistedLattice.covolume\_add\_index\_nonneg}, in \lean{grh/src/New/Reduction/Lattices/TwistedLattice.lean}.  In \lean{LongForms.lean} the proof of~(b) uses \lean{New.exists\_isRegularFamily} and \lean{New.Data.heightLe\_liftDeterminant\_ratio}, and Problems~\ref{pr:completion} and~\ref{pr:volume} are \lean{New.Data.mul\_completion\_cost\_le} and \lean{New.Data.heightLe\_chainDeterminant\_one\_div}.

\section*{Hints}

\subsection*{\S\ref{sec:toy}}
\begin{description}
\item[\ref{pr:toy-int}] \((\langle y_i, y_j \rangle) = N\, Y Y^{\mathsf T}\) for the matrix \(Y\) with rows \(y_i\), and \(\det Y\) is a nonzero integer.
\item[\ref{pr:toy-index}] Problem~\ref{pr:toy-int} for the rows of \(uY\).
\item[\ref{pr:toy-complete}] Fischer's inequality for the Gram matrix of \((y, z)\), Hadamard's for \(\Gram(z)\), then Problem~\ref{pr:toy-int}.
\item[\ref{pr:toy-steinitz}] Steinitz exchange and Problem~\ref{pr:toy-complete}; each standard basis vector has squared length \(N\).
\end{description}

\subsection*{\S\ref{sec:thmC}}
\begin{description}

\item[\ref{pr:scales}] \(m^2 \le 4x\), \(t \le 15x^5\) and \(x \ge 2\).
\item[\ref{pr:completion}] \(n' \le 4d_T/b^2\).
\item[\ref{pr:denom}] Input~(d), Problem~\ref{sd-pr:height-mul} and Problem~\ref{pr:scales}.
\item[\ref{pr:volume}] Invert~(c); bound the scalar by~(d) and the single lifts by~(b), and add by Problem~\ref{sd-pr:height-mul}.
\item[\ref{pr:C}] Apply~\eqref{eq:twist} with \(c_{\mathrm{short}} = 2x_1^2b^2\), \(c_{\mathrm{den}} = b\) and \(c_{\mathrm{vol}} = d_T(b + F - \main)\).
\end{description}

{\small
\begin{thebibliography}{GL}
\bibitem[GL]{GL} L.~Sernau, \emph{GRH for Dirichlet \(L\)-functions, in Lean~4}, a shorter rewriting of the Lean proof, text by Claude, \url{https://github.com/lsernau/grh-lean}, commit \texttt{72fead1}, readable edition \texttt{grh/}.
\bibitem[HJ]{HJ} R.~A.~Horn and C.~R.~Johnson, \emph{Matrix Analysis}, 2nd ed., Cambridge Univ.\ Press, 2013.
\end{thebibliography}
}

\end{document}
